\documentclass[12pt,a4paper]{article}

% Idioma y codificación
\usepackage[utf8]{inputenc}
\usepackage[T1]{fontenc}
\usepackage[spanish]{babel}
\usepackage{geometry}
\geometry{margin=1in}

% Espaciado entre párrafos
\usepackage{parskip}

% Listas, tablas y enlaces
\usepackage{enumitem}
\usepackage{booktabs}
\usepackage{hyperref}
\hypersetup{colorlinks=true, urlcolor=blue, linkcolor=black, citecolor=black}

% Estilo de secciones
\usepackage{titlesec}
\titleformat{\section}{\large\bfseries}{\thesection}{1em}{}
\titleformat{\subsection}{\normalsize\bfseries}{\thesubsection}{1em}{}
\titleformat{\subsubsection}{\normalsize\bfseries}{\thesubsubsection}{1em}{}
\setlist[itemize]{topsep=2pt,parsep=2pt,partopsep=0pt,leftmargin=1.5em}
\setlist[enumerate]{topsep=2pt,parsep=2pt,partopsep=0pt,leftmargin=1.8em}

% Datos del curso
\title{\textbf{FACULTAD DE ECONOMÍA Y NEGOCIOS \\ Departamento de Economía} \\
Big Data y Machine Learning para Economía Aplicada}
\author{Prof. Eduard Fernando Martínez González, Ph.D.}
\date{}

% Marca de agua en la primera página
\usepackage{eso-pic}
\usepackage{graphicx}
\AddToShipoutPictureBG*{
\AtPageUpperLeft{\raisebox{-3cm}{\hspace{2cm}\includegraphics[width=5cm]{pics/logo.png}}}}

%%=================================================
\begin{document}

\maketitle
\begin{itemize}
\item \textbf{Código-Materia:} 60-121 (NRC 10-993)
\item \textbf{Programa:} Maestría en Economía (ME)
\item \textbf{Prerrequisito:} Econometría I y II (o equivalente). Se recomienda manejo básico de \texttt{R}; no es requisito.
\item \textbf{Período académico:} 2026-2
\item \textbf{Horario y lugar:} martes de 8:00 a 11:00 a.\,m., salón 406-E (Bloque E)
\item \textbf{Fechas:} del 15 de septiembre al 24 de noviembre de 2026; nueve sesiones. No hay clase el 10 ni el 17 de noviembre.
\item \textbf{Intensidad semanal:} 3 horas
\item \textbf{Créditos:} 3
\item \textbf{Materiales:} \href{https://eduard-martinez.github.io/teaching/bdml-applied-economics/}{página del curso} y \href{https://github.com/eduard-martinez/bdml-applied-economics}{repositorio en GitHub} (presentaciones, prácticas de R con sus datos y talleres)
\item \textbf{Contacto:} \href{mailto:efmartinez@icesi.edu.co}{efmartinez@icesi.edu.co}
\end{itemize}

%%=================================================
\section{Descripción del curso}

En Inferencia Causal el objetivo es identificar correctamente un parámetro: el efecto $\beta$ de una intervención sobre un resultado. En este curso el objetivo cambia de lugar: nos interesa $Y$, es decir, construir modelos capaces de \textbf{predecir bien la variable de interés} con datos que el modelo no ha visto, saber \textbf{evaluar} esa capacidad predictiva con honestidad y \textbf{abrir el modelo} para entender de qué depende la predicción. Esa distinción entre explicar y predecir organiza todo el semestre y se marca desde la primera sesión.

El curso introduce a los estudiantes de la Maestría en Economía al marco del aprendizaje estadístico y a las herramientas de \emph{machine learning} (ML) más usadas hoy en economía aplicada: la regresión como máquina de predecir, la validación cruzada y la regularización (Ridge, Lasso, Elastic Net), la clasificación, los árboles, los \emph{random forests} y el \emph{boosting}, las herramientas para abrir la caja negra (importancia de variables, dependencia parcial, valores de Shapley) y las redes neuronales, con un puente final hacia el ML al servicio de la inferencia causal. El texto guía es \emph{An Introduction to Statistical Learning} (James, Witten, Hastie y Tibshirani, 2.ª ed.; ISL): cada sesión indica los capítulos y secciones que le corresponden.

El tratamiento combina \textbf{formalidad suficiente} con \textbf{mucha intuición}: para cada método se presenta el problema económico que resuelve, la ecuación que explica el mecanismo, una figura construida con los datos del curso y la implicación para la predicción; las derivaciones extensas viven en el apéndice de cada presentación. \textbf{Un solo caso empírico atraviesa las sesiones 2 a 8}: los puestos de votación de Colombia en la segunda vuelta presidencial de 2022, con sus características del censo de 2018 (los datos del \emph{working paper} de Gelvez, Cardiles, Martínez-González y Muñoz, 2026). Cada semana la teoría usa como ejemplos los números de ese caso, y la práctica de R que se publica con la sesión los reproduce todos. Los \emph{problem sets} usan una segunda base, avisos de venta de vivienda en Cali, y el proyecto final se hace con una base elegida por el estudiante.

\section{Resultados de aprendizaje}

Al finalizar el curso, el estudiante será capaz de:

\begin{itemize}
    \item Distinguir un problema de predicción de un problema de inferencia causal, y reconocer cuándo la capacidad predictiva fuera de muestra es la cantidad de interés.
    \item Construir el flujo de trabajo completo de un ejercicio predictivo: partición de los datos, preprocesamiento dentro del \emph{pipeline}, validación cruzada, ajuste de hiperparámetros, evaluación e interpretación.
    \item Evaluar modelos de regresión y clasificación con las métricas apropiadas (RMSE, MAE, $R^2$ fuera de muestra; matriz de confusión, AUC, calibración), con su incertidumbre y frente a un modelo de referencia, y diagnosticar dónde se concentran los errores.
    \item Aplicar y comparar regularización, árboles, \emph{random forests}, \emph{boosting} y redes neuronales sobre datos económicos, sabiendo cuándo conviene cada familia.
    \item Interpretar modelos flexibles con importancia de variables, dependencia parcial y valores SHAP, distinguiendo lo que describen (de qué depende la predicción) de lo que no establecen (efectos causales).
    \item Leer críticamente y presentar artículos de economía que usan ML, evaluando la pertinencia del método, de la validación y de la interpretación.
    \item Implementar todo lo anterior en \texttt{R} con código transparente y sin cajas negras de programación (los paquetes del curso: \texttt{glmnet}, \texttt{rpart}, \texttt{randomForest}, \texttt{xgboost} y \texttt{nnet}): las prácticas escriben cada paso.
\end{itemize}

%%=================================================
\section{Contenido}

Las nueve sesiones arrancan en el marco del aprendizaje estadístico, construyen la regresión como \emph{baseline} y el par validación cruzada--regularización como el \emph{pipeline} honesto, pasan a la clasificación, los ensambles de árboles, la interpretación y la representación (PCA, \emph{clustering} y redes), y cierran con la sesión de \emph{papers} y síntesis y con la presentación del proyecto final. El Problem Set 1 cubre las sesiones 2 a 4 y el Problem Set 2 las sesiones 5 a 7.

\subsection{Calendario de sesiones y entregas}

Las clases son los martes de 8:00 a 11:00 a.\,m. La sesión 9 cae dentro del período de exámenes finales del calendario académico 2026-2 (17 al 28 de noviembre); las notas definitivas se registran a partir del 1 de diciembre.

\begin{center}
\small
\begin{tabular}{@{}clp{6.2cm}p{4.7cm}@{}}
\toprule
S & Fecha & Tema (secciones de ISL) & Hitos \\
\midrule
1 & mar 15 sep & Predecir, no explicar: el marco del aprendizaje estadístico (caps.\ 1--2) & \\
2 & mar 22 sep & La regresión lineal como máquina de predecir (2.2, 3 y 3.5; cap.\ 7 en una lámina) & Se eligen las áreas de los \emph{papers} de la sesión 8 \\
3 & mar 29 sep & Validación cruzada y regularización: el juez y la perilla (5.1, 6.2 y 6.4) & \\
4 & mar 6 oct & Clasificación: de la probabilidad a la decisión (4.1--4.3 y 4.7.6) & Se publica el Problem Set 1 \\
5 & mar 13 oct & Árboles, bosques y \emph{boosting}: partir, promediar y corregir (cap.\ 8) & \\
6 & mar 20 oct & Abrir la caja negra: qué aprendió la máquina (Molnar, caps.\ seleccionados) & Entrega del PS1 (8:00 a.\,m.); se publica el PS2 \\
7 & mar 27 oct & Representar: PCA, $k$-medias y redes neuronales (12.2 y 12.4.1; 10.1--10.3 y 10.6--10.7) & \emph{Debrief} del \emph{ranking} del PS1 \\
8 & mar 3 nov & Los \emph{papers} y la síntesis del caso; puente causal con DML & Presentaciones de \emph{papers}; Entrega 1 del proyecto (idea + \emph{pitch}) \\
--- & mar 10 nov & Sin clase (asesoría opcional del proyecto) & Entrega del PS2 (11:59 p.\,m.) \\
--- & mar 17 nov & Sin clase & \\
9 & mar 24 nov & Presentación del proyecto final & Entrega 2 del proyecto (documento); \emph{debrief} del PS2 \\
\bottomrule
\end{tabular}
\end{center}

\subsection{Módulo 1: El marco, el \emph{baseline} y el \emph{pipeline} honesto}

\subsubsection{Sesión 1 --- Predecir, no explicar: el marco del aprendizaje estadístico}
\textbf{Temas:} de $\beta$ a $Y$ (los mismos datos, dos preguntas; \emph{prediction policy problems}); qué es aprender $f$; flexibilidad vs.\ interpretabilidad; error de entrenamiento vs.\ error de prueba, sobreajuste y la curva U; sesgo--varianza con intuición; las tareas (regresión y clasificación, supervisado y no supervisado) y el flujo de trabajo del curso. \\
\textbf{Lecturas:} ISL, caps.\ 1--2; Mullainathan \& Spiess (2017); Kleinberg, Ludwig, Mullainathan \& Obermeyer (2015). \\
\textbf{Aplicación en R:} un mismo conjunto de datos, tres modelos: qué pasa con cada uno cuando llega un dato nuevo.

\subsubsection{Sesión 2 --- La regresión lineal como máquina de predecir}
\textbf{Temas:} describir el \emph{target} antes de modelar; la regresión simple (el RSS, la solución como proyección, las métricas en un puesto concreto); más flexible que la recta con una variable (polinomios, escalones y \emph{splines}, $k$-NN y la maldición de la dimensionalidad); la regresión múltiple (mirar antes de ajustar, el plano, RSE y $R^2$ releídos, la gemela y la colinealidad); después de predecir (predicho vs.\ observado, residuos como detector, la referencia y la tabla canónica); extrapolación y la prueba de fuego: la recta no viaja. \\
\textbf{Lecturas:} ISL, secciones 2.2 y 3.5 y cap.\ 3 (cap.\ 7 en una lámina). \\
\textbf{Aplicación en R:} Cali y su área metropolitana: la media como referencia, la recta, $k$-NN en el mapa y las 511 combinaciones de nueve covariables; la selección con la prueba queda en evidencia (el gancho de la sesión 3).

\subsubsection{Sesión 3 --- Validación cruzada y regularización: el juez y la perilla}
\textbf{Temas:} \emph{el juez}: una partición de validación es una lotería; la validación cruzada con $J$ pliegues y el error estándar del promedio; el protocolo del curso (el mínimo y la regla de un error estándar); el concurso de los 511 con juez honesto. \emph{La perilla}: cobrar por la complejidad; Ridge divide, el Lasso resta y apaga; Elastic Net; los caminos de coeficientes y la elección de $\lambda$ por validación cruzada. \emph{En el país}: las 1.291 columnas de censo $\times$ departamento, OLS se desploma, el juez elige al Lasso; dos jueces, dos preguntas (pliegues al azar o municipios enteros). \\
\textbf{Lecturas:} ISL, secciones 5.1, 6.2 y 6.4; Zou \& Hastie (2005) para Elastic Net. \\
\textbf{Aplicación en R:} las tres partes de la teoría sobre los 12.000 puestos del país, y un segundo script corto que hace \emph{zoom} al juez (el error de cada pliegue y su EE) y a los caminos de Ridge y Lasso.

\subsubsection{Sesión 4 --- Clasificación: de la probabilidad a la decisión}
\textbf{Temas:} la $y$ binaria (gana o no gana el puesto) y su distribución; las referencias sin modelo; del modelo de probabilidad lineal (LPM) a la logística (verosimilitud, log-loss, cómo leer los coeficientes y qué no leer); la matriz de confusión y sus métricas (exactitud, precisión, sensibilidad o \emph{recall} y especificidad), por qué la exactitud engaña con clases desbalanceadas y qué mirar cuando los positivos son raros (precisión--\emph{recall}); el umbral como decisión económica ($\tau^* = c_{FP}/(c_{FP}+c_{FN})$) y la curva de costo; ROC y AUC (ordenar sin fijar umbral), calibración y Brier; los competidores con la perilla de la sesión 3 (LPM con Lasso, la logística con las variables seleccionadas, el Lasso logístico) y los vecinos del mapa. \\
\textbf{Lecturas:} ISL, secciones 4.1--4.3 (clasificación y regresión logística) y 4.7.6 (el laboratorio de $k$-NN como clasificador); de la sección 4.4.2, la matriz de confusión y la curva ROC. LDA, QDA y Naive Bayes (secciones 4.4--4.5) quedan en el apéndice de la sesión. Fawcett (2006). \\
\textbf{Aplicación en R:} ¿gana Petro el puesto? Referencias, LPM vs.\ logística, matriz de confusión y métricas con el desbalance de 2010, umbral con costos, ROC/AUC, calibración y la tabla final de clasificadores.

\subsection{Módulo 2: Árboles, ensambles e interpretación}

\subsubsection{Sesión 5 --- Árboles, bosques y \emph{boosting}: partir, promediar y corregir}
\textbf{Temas:} el árbol (el primer corte a mano, la partición binaria recursiva, interacciones sin escribirlas, poda por costo--complejidad, la varianza del árbol); el bosque (\emph{bagging}, la correlación entre árboles, el error \emph{out-of-bag}, la perilla $m$, por qué no extrapola); el \emph{boosting} (ajustar residuos; las perillas $B$, $\eta$ y $d$ con el juez de la sesión 3); la carrera del caso: cuánto valen las interacciones, el mapa como predictor y el promedio de modelos. \\
\textbf{Lecturas:} ISL, cap.\ 8 (8.1, 8.2.1--8.2.3 y 8.2.5); Breiman (2001). \\
\textbf{Aplicación en R:} del corte a mano en diez puestos de Cali a la carrera nacional (\texttt{rpart}, \texttt{randomForest}, \texttt{xgboost}) en las dos preguntas del caso.

\subsubsection{Sesión 6 --- Abrir la caja negra: qué aprendió la máquina}
\textbf{Temas:} la importancia (ganancia, permutación, ablación por bloques y sus trampas); la dependencia (PDP, ICE y ALE; cuando el PDP pregunta por datos que no existen); valores de Shapley y SHAP (el reparto justo, TreeSHAP, cascadas de puestos concretos); el efecto Rashomon cuando entra el mapa; a quién le falla el modelo (subgrupos, extremos y equidad). Hoy no entra ningún modelo nuevo: se abre el \emph{boosting} de la sesión 5. \\
\textbf{Lecturas:} Molnar (2025), capítulos de importancia por permutación, dependencia parcial y SHAP; Lundberg \& Lee (2017). \\
\textbf{Aplicación en R:} las tres importancias lado a lado, PDP/ICE/ALE, Shapley a mano en un puesto de Quibdó, SHAP con \texttt{predcontrib} y el error por subgrupos.

\subsubsection{Sesión 7 --- Representar: PCA, $k$-medias y redes neuronales}
\textbf{Temas:} comprimir con sentido: PCA (la dirección de máxima varianza, cuántas componentes, las cargas como índice, PCR con el juez); $k$-medias (el codo y la silueta, las ``cuatro Colombias'', ¿los grupos predicen?); la red neuronal (de la componente principal a la unidad oculta, por qué hace falta la no linealidad, capacidad y \emph{weight decay} con el juez, descenso de gradiente y parada temprana, el papel de la semilla); convolución y \emph{transfer learning} con imágenes. Abre con el \emph{debrief} del PS1. \\
\textbf{Lecturas:} ISL, secciones 12.2 y 12.4.1 (PCA y $k$-medias) y 10.1--10.3 y 10.6--10.7 (redes); Jean et al.\ (2016). \\
\textbf{Aplicación en R:} el PCA de las 37 variables del censo, las cuatro Colombias, la grilla de redes con el juez y la carrera contra el \emph{boosting}; el voto como imagen y un filtro de bordes.

\subsection{Módulo 3: Síntesis, papers y proyecto}

\subsubsection{Sesión 8 --- Los \emph{papers} y la síntesis del caso}
\textbf{Temas:} seis presentaciones de \emph{papers} por áreas de la economía (15 minutos cada una); la síntesis del caso (la cascada del curso: cuánto aportaron la información, la flexibilidad y la combinación de modelos; la tabla ingredientes $\times$ familias; ¿la estructura viaja a municipios nuevos, a otros grupos, a otro año?); los \emph{pitches} de la Entrega 1 del proyecto; y el puente causal: ML al servicio de $\beta$ con \emph{double machine learning} sobre los mismos datos del curso. \\
\textbf{Lecturas:} Athey \& Imbens (2019), secciones 1--3 (preparan el puente causal), y el \emph{paper} elegido de su área; Chernozhukov et al.\ (2018) como referencia del DML. \\
\textbf{Hitos:} presentaciones de \emph{papers} (10\% de la nota) y Entrega 1 del proyecto (idea en una página + \emph{pitch}).

\subsubsection{Sesión 9 --- Presentación del proyecto final}
\textbf{Temas:} cada estudiante o pareja presenta su propuesta de investigación predictiva (20 minutos más preguntas); \emph{debrief} del \emph{ranking} del PS2; lo que se llevan del curso, lo que no vimos y dónde seguir. \\
\textbf{Hitos:} Entrega 2 del proyecto (documento) y presentación.

%%=================================================
\section{Metodología}

Las sesiones combinan tres piezas:

\begin{enumerate}
    \item \textbf{Exposición teórica:} cada concepto entra por una necesidad del caso y sigue la misma secuencia: pregunta $\rightarrow$ intuición $\rightarrow$ formalización $\rightarrow$ figura $\rightarrow$ implicación para la predicción. Los ejemplos numéricos de la teoría son los del caso de la semana: \textbf{en la teoría no seguimos un paper; seguimos el ejemplo que después se replica en código}. Las derivaciones extensas viven en el apéndice ``Para profundizar'' de cada presentación.

    \item \textbf{Práctica de R:} cada sesión publica en el repositorio un script plano y comentado (con sus datos en \texttt{input/} y un \texttt{.zip} listo para descargar) que reproduce \textbf{todas} las cifras y figuras de la presentación. El código evita atajos opacos: las grillas se llenan con un \texttt{for}, el RMSE se escribe en línea y las predicciones se comparan siempre contra el \emph{target} observado.

    \item \textbf{El caso y los \emph{papers}:} el hilo empírico de las sesiones 2 a 8 son los puestos de votación de 2022 con el censo de 2018: la pregunta de regresión (el porcentaje de voto) y la de clasificación (¿gana el puesto?) se heredan y mejoran semana a semana, y los mejores modelos del curso tienen memoria. La literatura se organiza \textbf{por áreas de la economía} (economía política; crimen y justicia; salud; desarrollo y pobreza; trabajo; vivienda; macroeconomía y pronóstico; finanzas y crédito; organización industrial): en la sesión 2 cada estudiante elige un área, y en la sesión 8 presenta uno de sus dos \emph{papers} candidatos (carpeta \texttt{literature/} del repositorio). Los artículos con carreras de caballos de algoritmos (Gu, Kelly \& Xiu, 2020; Medeiros et al., 2021; Goulet Coulombe et al., 2022) son lecturas de referencia para el proyecto.
\end{enumerate}

Cuatro reglas de oro se repiten durante el semestre: (1) el error de entrenamiento es optimista; (2) la prueba se abre una sola vez; (3) la validación cruzada valida \emph{pipelines} completos, no modelos sueltos; y (4) la referencia es lo que la decisión haría sin el modelo.

\section{Otras obligaciones de los estudiantes}

Cada estudiante debe realizar de manera autónoma una nivelación en \texttt{R} (introducción al lenguaje, manipulación de datos con \texttt{dplyr} y gráficos básicos) antes de las sesiones aplicadas; se proporciona material de referencia. Las prácticas publicadas de cada semana son, además, el material de estudio aplicado: se espera que el estudiante las corra y las entienda línea a línea antes de la sesión siguiente.

%%=================================================
\section{Evaluación}

La calificación final del curso se basa en cuatro componentes:

\begin{enumerate}
    \item \textbf{Presentación de \emph{paper} en la sesión 8 (10\%)} \\
    Cada estudiante elige un área de la economía en la sesión 2 y presenta en la sesión 8, en 15 minutos, uno de los dos \emph{papers} candidatos de su área. El guion: la pregunta y la decisión; los datos y la unidad de análisis; el modelo de referencia; la evaluación fuera de muestra y dónde falla; el resultado principal; y qué método del curso es. No hay guías de lectura semanales: los \emph{papers} del semestre se conversan en clase como ejemplos.

    \item \textbf{Problem Set 1: referencia, regresión, validación y clasificación (15\%)} \\
    En parejas, sobre una única base de avisos de venta de vivienda en Cali con diccionario y partición fija; el precio y la condición de vivienda VIS de los avisos de prueba están ocultos. Partes y puntos: (1) referencia y regresión (sesión 2, 15 pts); (2) validación y regularización (sesión 3, 25 pts); (3) clasificación (sesión 4, 15 pts); (4) modelo final, con la prueba abierta una sola vez (15 pts); \emph{ranking} del curso por el RMSE de prueba del modelo final (20 pts: el primero 20, el último 10); y reproducibilidad --- \texttt{00\_run.R} corre de principio a fin (10 pts). Se publica el martes 6 de octubre (sesión 4) y se entrega el martes 20 de octubre a las 8:00 a.\,m.; \emph{debrief} del \emph{ranking} en la sesión 7.

    \item \textbf{Problem Set 2: árboles, caja abierta, redes y el modelo final del curso (15\%)} \\
    Misma base y mismas reglas. Partes y puntos: (1) árboles, bosques y \emph{boosting} (sesión 5, 20 pts); (2) abrir la caja (sesión 6, 25 pts); (3) la red contra el bosque (sesión 7, 10 pts); (4) modelo final del curso y lectura de una página (15 pts); \emph{ranking} (20 pts) y reproducibilidad (10 pts). Se publica el martes 20 de octubre (sesión 6) y se entrega el martes 10 de noviembre a las 11:59 p.\,m.; \emph{debrief} del \emph{ranking} el 24 de noviembre.

    \item \textbf{Proyecto final: una propuesta de investigación predictiva (60\%)} \\
    El producto del curso no es un informe con resultados: es una \textbf{propuesta de investigación predictiva} sobre una base elegida por el estudiante, con la pregunta y la decisión que la motiva, los datos, la métrica y el modelo de referencia, la estrategia de validación, los métodos que se usarían y por qué, los riesgos y el plan de trabajo. Sobre 100 puntos del componente:
    \begin{itemize}
        \item \textbf{Entrega 1 --- La idea y el \emph{pitch} (10 pts), 3 de noviembre:} la idea en una página y un \emph{pitch} corto en la sesión 8.
        \item \textbf{Entrega 2 --- El documento (50 pts), 24 de noviembre:} máximo 8 páginas.
        \item \textbf{Presentación (40 pts), 24 de noviembre:} 20 minutos más preguntas, en la sesión 9.
    \end{itemize}
    Si el proyecto se hace individual o en parejas se confirma en clase. Los lineamientos detallados se publican en la carpeta \texttt{final-project/} del repositorio.
\end{enumerate}

\noindent \textbf{Resumen de la evaluación:}
\begin{itemize}
    \item Presentación de \emph{paper} (sesión 8): 10\%
    \item Problem Set 1 (sesiones 2--4): 15\%
    \item Problem Set 2 (sesiones 5--7): 15\%
    \item Proyecto final (idea y \emph{pitch}, documento y presentación): 60\%
\end{itemize}

\noindent \textbf{Nota sobre el uso de IA:} Se permite el uso de herramientas de IA en cualquier componente del curso, siempre que se declare explícitamente qué se usó y para qué.

%%=================================================
\section{Bibliografía}

\subsection*{Textos base}
\begin{itemize}
    \item James, G., Witten, D., Hastie, T., \& Tibshirani, R.\ (2021). \textit{An Introduction to Statistical Learning with Applications in R} (2.ª ed.). Springer. [ISL, el texto guía] \href{https://www.statlearning.com}{[disponible online]}
    \item Hastie, T., Tibshirani, R., \& Friedman, J.\ (2009). \textit{The Elements of Statistical Learning} (2.ª ed.). Springer. [ESL]
    \item Molnar, C.\ (2025). \textit{Interpretable Machine Learning: A Guide for Making Black Box Models Explainable} (3.ª ed.). \href{https://christophm.github.io/interpretable-ml-book/}{[disponible online]}
\end{itemize}

\subsection*{Marco: ML en economía}
\begin{itemize}
    \item Mullainathan, S., \& Spiess, J.\ (2017). Machine Learning: An Applied Econometric Approach. \textit{Journal of Economic Perspectives}, 31(2), 87--106.
    \item Kleinberg, J., Ludwig, J., Mullainathan, S., \& Obermeyer, Z.\ (2015). Prediction Policy Problems. \textit{American Economic Review: Papers \& Proceedings}, 105(5), 491--495.
    \item Athey, S., \& Imbens, G.\ W.\ (2019). Machine Learning Methods That Economists Should Know About. \textit{Annual Review of Economics}, 11, 685--725.
    \item Varian, H.\ R.\ (2014). Big Data: New Tricks for Econometrics. \textit{Journal of Economic Perspectives}, 28(2), 3--28.
\end{itemize}

\subsection*{El caso del curso}
\begin{itemize}
    \item Gelvez, J.\ D., Cardiles, A., Martínez-González, E.\ F., \& Muñoz, M.\ (2026). How Predictable Is an Election? A Machine-Learning Approach to Electoral Behavior in Colombia. Documento de trabajo. [De aquí salen los datos de los puestos de votación.]
\end{itemize}

\subsection*{Literatura por áreas de la economía}

La carpeta \texttt{literature/} del repositorio organiza los artículos aplicados por área, con dos candidatos por área para las presentaciones de la sesión 8. Una selección:

\begin{itemize}
    \item \textbf{Economía política:} Kim \& Zilinsky (2024), Division Does Not Imply Predictability, \textit{Political Behavior}, 46, 67--87 · Gentzkow, Shapiro \& Taddy (2019), Measuring Group Differences in High-Dimensional Choices, \textit{Econometrica}, 87(4).
    \item \textbf{Crimen y justicia:} Kleinberg, Lakkaraju, Leskovec, Ludwig \& Mullainathan (2018), Human Decisions and Machine Predictions, \textit{Quarterly Journal of Economics}, 133(1) · Angelova, Dobbie \& Yang (2023), Algorithmic Recommendations and Human Discretion, documento de trabajo · Ludwig \& Mullainathan (2024), Machine Learning as a Tool for Hypothesis Generation, \textit{Quarterly Journal of Economics}, 139(2).
    \item \textbf{Salud:} Mullainathan \& Obermeyer (2022), Diagnosing Physician Error, \textit{Quarterly Journal of Economics}, 137(2) · Einav, Finkelstein, Mullainathan \& Obermeyer (2018), Predictive Modeling of U.S.\ Health Care Spending in Late Life, \textit{Science}, 360(6396).
    \item \textbf{Desarrollo y pobreza:} Aiken, Bedoya, Blumenstock \& Coville (2023), Program Targeting with Machine Learning and Mobile Phone Data, \textit{Journal of Development Economics}, 161 · Haushofer, Niehaus, Paramo, Miguel \& Walker (2025), Targeting Impact versus Deprivation, documento de trabajo · Blumenstock, Cadamuro \& On (2015), Predicting Poverty and Wealth from Mobile Phone Metadata, \textit{Science}, 350(6264).
    \item \textbf{Trabajo y capital humano:} Hoffman, Kahn \& Li (2018), Discretion in Hiring, \textit{Quarterly Journal of Economics}, 133(2) · Chalfin et al.\ (2016), Productivity and Selection of Human Capital with Machine Learning, \textit{American Economic Review: P\&P}, 106(5).
    \item \textbf{Vivienda y economía urbana} (la base de los \emph{problem sets}): Bogin \& Shui (2020), Appraisal Accuracy and Automated Valuation Models in Rural Areas, \textit{J.\ of Real Estate Finance and Economics}, 60 · Deppner \& Cajias (2024), Accounting for Spatial Autocorrelation in Algorithm-Driven Hedonic Models, \textit{J.\ of Real Estate Finance and Economics}, 68.
    \item \textbf{Macroeconomía y pronóstico:} Medeiros, Vasconcelos, Veiga \& Zilberman (2021), Forecasting Inflation in a Data-Rich Environment, \textit{J.\ of Business \& Economic Statistics}, 39(1) · Goulet Coulombe, Leroux, Stevanovic \& Surprenant (2022), How Is Machine Learning Useful for Macroeconomic Forecasting?, \textit{J.\ of Applied Econometrics}, 37(5) · Bluwstein et al.\ (2023), Credit Growth, the Yield Curve and Financial Crisis Prediction, \textit{J.\ of International Economics}, 145.
    \item \textbf{Finanzas y crédito:} Fuster, Goldsmith-Pinkham, Ramadorai \& Walther (2022), Predictably Unequal? The Effects of Machine Learning on Credit Markets, \textit{Journal of Finance}, 77(1) · Dobbie, Liberman, Paravisini \& Pathania (2021), Measuring Bias in Consumer Lending, \textit{Review of Economic Studies}, 88(6) · Gu, Kelly \& Xiu (2020), Empirical Asset Pricing via Machine Learning, \textit{Review of Financial Studies}, 33(5).
    \item \textbf{Organización industrial:} Dubé \& Misra (2023), Personalized Pricing and Consumer Welfare, \textit{Journal of Political Economy}, 131(1) · Bandiera, Prat, Hansen \& Sadun (2020), CEO Behavior and Firm Performance, \textit{Journal of Political Economy}, 128(4).
\end{itemize}

\subsection*{Metodológicos y de método}
\begin{itemize}
    \item Zou, H., \& Hastie, T.\ (2005). Regularization and Variable Selection via the Elastic Net. \textit{Journal of the Royal Statistical Society: Series B}, 67(2), 301--320.
    \item Fawcett, T.\ (2006). An Introduction to ROC Analysis. \textit{Pattern Recognition Letters}, 27(8), 861--874.
    \item Breiman, L.\ (2001). Random Forests. \textit{Machine Learning}, 45(1), 5--32.
    \item Chen, T., \& Guestrin, C.\ (2016). XGBoost: A Scalable Tree Boosting System. \textit{Proceedings of KDD '16}, 785--794.
    \item Lundberg, S.\ M., \& Lee, S.-I.\ (2017). A Unified Approach to Interpreting Model Predictions. \textit{Advances in Neural Information Processing Systems}, 30.
    \item Friedman, J., Hastie, T., \& Tibshirani, R.\ (2010). Regularization Paths for Generalized Linear Models via Coordinate Descent. \textit{Journal of Statistical Software}, 33(1), 1--22. [\texttt{glmnet}]
    \item Jean, N., Burke, M., Xie, M., Davis, W.\ M., Lobell, D.\ B., \& Ermon, S.\ (2016). Combining Satellite Imagery and Machine Learning to Predict Poverty. \textit{Science}, 353(6301), 790--794.
    \item Filmer, D., \& Pritchett, L.\ (2001). Estimating Wealth Effects Without Expenditure Data --- or Tears. \textit{Demography}, 38(1), 115--132.
    \item Belloni, A., Chernozhukov, V., \& Hansen, C.\ (2014). High-Dimensional Methods and Inference on Structural and Treatment Effects. \textit{Journal of Economic Perspectives}, 28(2), 29--50.
    \item Chernozhukov, V., Chetverikov, D., Demirer, M., Duflo, E., Hansen, C., Newey, W., \& Robins, J.\ (2018). Double/Debiased Machine Learning for Treatment and Structural Parameters. \textit{Econometrics Journal}, 21(1), C1--C68.
\end{itemize}

\end{document}
